\begingroup
\renewcommand{\arraystretch}{1.45}
\fancyhead[L]{\small DETERMINATION OF ATTRACTION}
\Needspace{44\baselineskip}
\AttractionArticle{16}

As for the quantities $P,Q$, each evidently becomes
\[
\begin{aligned}
&=\frac{1}{2(G+G')^{\frac32}}
\end{aligned}
\]
whenever $G'=G''$, whereas in all other cases they must be referred to transcendental quantities. How these may be expressed by series is well known. We hope, however, that readers will welcome our taking this opportunity to explain the determination of these and other transcendental quantities by a special, most convenient algorithm, which we have frequently used for many years now, and which we intend to discuss more fully elsewhere.

\Needspace{36\baselineskip}
Let $m,n$ be two positive quantities, and let us set
\[
\begin{aligned}
&m'=\frac12(m+n),\qquad n'=\surd{mn}
\end{aligned}
\]
so that $m',n'$ are respectively the arithmetic and geometric means between $m$ and $n$. We shall assume that the geometric mean is always taken as positive. Similarly, let
\[
\begin{aligned}
&m''=\frac12(m'+n'),\qquad n''=\surd{m'n'}\\[8pt]
&m'''=\frac12(m''+n''),\qquad n'''=\surd{m''n''}
\end{aligned}
\]
and so on. In this way the sequences $m,m',m'',m'''$ etc., and $n,n',n'',n'''$ etc. will converge very rapidly towards a \emph{common limit}, which we shall denote by $\mu$ and simply call the \emph{arithmetic-geometric mean} between $m$ and $n$. We shall now prove that $\frac1\mu$ is the value of the integral
\[
\begin{aligned}
&\int\frac{\mathrm dT}{2\pi\surd{(mm\cos T^2+nn\sin T^2)}}
\end{aligned}
\]
taken from $T=0$ to $T=360^0$.

\Needspace{40\baselineskip}
\emph{Proof.} Let us suppose that the variable $T$ is expressed in terms of another $T'$ so that
\[
\begin{aligned}
&\sin T=\frac{2m\sin T'}{(m+n)\cos T^{\prime 2}+2m\sin T^{\prime 2}}
\end{aligned}
\]
and it will readily be seen that, as $T'$ increases from the value $0$ to $90^0,180^0,270^0,360^0$, $T$ too increases, although at unequal intervals, from $0$ to $90^0,180^0,270^0,360^0$. On carrying out the expansion properly, it is found that
\[
\begin{aligned}
&\frac{\mathrm dT}{\surd{(mm\cos T^2+nn\sin T^2)}}
=\frac{\mathrm dT'}{\surd{(m'm'\cos T^{\prime 2}+n'n'\sin T^{\prime 2})}}
\end{aligned}
\]
\SourceLeaf{362}
\Needspace{34\baselineskip}
and hence the values of the integrals
\[
\begin{gathered}\int\frac{\mathrm{d}T}{2\pi\surd(mm\cos T^2+nn\sin T^2)},\\[7pt]
\int\frac{\mathrm{d}T'}{2\pi\surd(m'm'\cos T'^2+n'n'\sin T'^2)}\end{gathered}
\]
are equal when each variable is taken from the value $0$ to the value $360^{0}$. Since this process may be continued in the same way, it is clear that these values also equal the value of the integral
\[
\int\frac{\mathrm{d}\Theta}{2\pi\surd(\mu\mu\cos\Theta^2+\mu\mu\sin\Theta^2)}
\]
from $\Theta=0$ to $\Theta=360^{0}$, which evidently becomes $=\frac{1}{\mu}$. Q.~E.~D.

\Needspace{44\baselineskip}
\AttractionArticle{17}
From the equation expressing the relation between $T$ and $T'$,
\[
(m-n)\sin T.\sin T'^2=2m\sin T'-(m+n)\sin T
\]
it readily follows that
\[
\begin{aligned}
\surd(mm\cos T^2+nn\sin T^2)&=m-(m-n)\sin T.\sin T'\\[6pt]
\surd(m'm'\cos T'^2+n'n'\sin T'^2)&=m\mathop{\mathrm{cotang}}T.\mathop{\mathrm{tang}}T'
\end{aligned}
\]
and hence, with the aid of the same equation,
\[
\begin{aligned}
&\sin T.\sin T'.\surd(mm\cos T^2+nn\sin T^2)+m'(\cos T^2-\sin T^2)\\[6pt]
&=\cos T.\cos T'.\surd(m'm'\cos T'^2+n'n'\sin T'^2)-\tfrac12(m-n)\sin T'^2
\end{aligned}
\]
\par Multiplying this equation by
\[
\frac{\mathrm{d}T}{\surd(mm\cos T^2+nn\sin T^2)}
=\frac{\mathrm{d}T'}{\surd(m'm'\cos T'^2+n'n'\sin T'^2)}
\]
gives
\[
\begin{aligned}&\frac{m'(\cos T^2-\sin T^2)\mathrm{d}T}{\surd(mm\cos T^2+nn\sin T^2)}\\[7pt]&=-\frac{\tfrac12(m-n)\sin T'^2\mathrm{d}T'}{\surd(m'm'\cos T'^2+n'n'\sin T'^2)}+\mathrm{d}.\sin T'\cos T\end{aligned}
\]
\Needspace{30\baselineskip}
\par Multiplying this equation by $\frac{m-n}{\pi}$, substituting
$m'(m-n)=\tfrac12(mm-nn)$,
$(m-n)^2=4(m'm'-n'n')$,
$\sin T'^2=\tfrac12-\tfrac12(\cos T'^2-\sin T'^2)$, and integrating from the values $T$ and $T'=0$ to $360^{0}$, we have:
\[
\begin{aligned}
&(mm-nn)\int\frac{(\cos T^2-\sin T^2).\mathrm{d}T}{2\pi\surd(mm\cos T^2+nn\sin T^2)}\\[6pt]
&=-\frac{2(m'm'-n'n')}{\mu}\\[7pt]&\quad+2(m'm'-n'n')\int\frac{(\cos T'^2-\sin T'^2)\mathrm{d}T'}{2\pi\surd(m'm'\cos T'^2+n'n'\sin T'^2)}
\end{aligned}
\]
\SourceLeaf{363}
Since the definite integral on the right may be transformed in the same way, the integral
\[
\int\frac{(\cos T^2-\sin T^2)\mathrm{d}T}{2\pi\surd(mm\cos T^2+nn\sin T^2)}
\]
will evidently be expressed by the very rapidly convergent infinite series
\[
-\frac{\begin{gathered}2(m'm'-n'n')+4(m''m''-n''n'')\\+8(m'''m'''-n'''n''')+\mathrm{etc.}\end{gathered}}{(mm-nn)\mu}
=-\frac{\nu}{\mu}
\]
\Needspace{31\baselineskip}
\par The numerical calculation is most conveniently carried out by logarithms if we set
\[
\begin{gathered}\tfrac14\surd(mm-nn)=\lambda,\quad\tfrac14\surd(m'm'-n'n')=\lambda',\\[6pt]\tfrac14\surd(m''m''-n''n'')=\lambda''\ \mathrm{etc.}\end{gathered}
\]
whence it will follow that
\[
\lambda'=\frac{\lambda\lambda}{m'},\quad
\lambda''=\frac{\lambda'\lambda'}{m''},\quad
\lambda'''=\frac{\lambda''\lambda''}{m'''}\ \mathrm{etc.}\quad\text{and}
\]
\[
\nu=\frac{2\lambda'\lambda'+4\lambda''\lambda''+8\lambda'''\lambda'''+\mathrm{etc.}}{\lambda\lambda}
\]

\clearpage
\endgroup
